\begin{prerequis}
Développer, factoriser par un facteur commun, résoudre une équation du premier
degré et utiliser un tableau de signes simple.
\end{prerequis}
\begin{automatismes}
Développer $(x-3)^2$; factoriser $5x^2-10x$; résoudre $3x-7=2$;
donner le signe de $(x-1)(x+4)$.
\end{automatismes}
\begin{activite}[Activité d'entrée -- Deux dimensions, même aire]
Un rectangle a pour côtés $x+2$ et $x+5$. Pour quelles valeurs de $x$ son aire
vaut-elle $28$? Mettre le problème sous forme d'équation, chercher d'abord par
essais organisés, puis transformer l'expression pour faire apparaître un trinôme.
Comparer les stratégies et discuter les contraintes géométriques sur $x$.
\end{activite}

